\section{Results}\label{sec:Results}\justifying{}
\subsection{Database of enalapril pharmacokinetics}
A comprehensive literature search on the PKPD of enalapril was performed in the literature databases~\href{https://www.pkpdai.com/}{PKPDAI}~\cite{GonzalezHernandez2021} and ~\href{https://pubmed.ncbi.nlm.nih.gov/}{PubMed}. From over 600 trials assessed, 91 were prioritised according to the following criteria: (1) the study had to be conducted in humans; (2) the study had to be conducted in adult subjects; (3) the study had to be conducted \textit{in vivo}; (4) the study had to report time-course data on plasma or serum and urine concentrations of enalapril and/or its metabolite enalaprilat. Particular attention was paid to studies in subjects with hypertension and renal, hepatic and cardiac impairment. Based on these criteria, 49 studies were selected for data curation, which formed the database for the development and evaluation of the PBPK model. 
% Data was curated from Arafat et al. 2005~\cite{Arafat2005}, Baba et al. 1990~\cite{Baba1990}, Biollaz et al 1982~\cite{Biollaz1982}, Dayyih et al. 2017~\cite{Dayyih2017}, Dickstein et al. 1987~\cite{Dickstein1987}, Donnelly et al. 1990~\cite{Donnelly1990}, Fruncillo et al. 1987~\cite{Fruncillo1987}, Ghosh et al. 2012~\cite{Ghosh2012}, Gu et al. 2004~\cite{Gu2004}, He et al. 2014~\cite{He2014}, Her et al. 2021~\cite{Her2021}, Hockings et al. 1986~\cite{Hockings1986}, Howes et al. 1991~\cite{Howes1991}, Ishizaki et al. 1988~\cite{Ishizaki1988}, Johnston et al. 1992~\cite{Johnston1992}, Kelly et al. 1986~\cite{Kelly1986}, Lee et al. 2003~\cite{Lee2003}, Lees et al. 1987~\cite{Lees1987a}, Lowenthal et al. 1985~\cite{Lowenthal1985}, Lu et al. 2009~\cite{Lu2009}, MacDonald et al. 1993~\cite{MacDonald1993}, Marzo et al. 2002~\cite{Marzo2002}, Matalka et al. 2002~\cite{Matalka2002}, Moffett et al. 2014~\cite{Moffett2014}, Mujais et al. 1992~\cite{Mujais1992}, Najib et al. 2003~\cite{Najib2003a}, Niopas et al. 2004~\cite{Niopas2004}, Oguchi et al. 1993~\cite{Oguchi1993}, Ohnishi et al. 1989~\cite{Ohnishi1989}, Portoles et al. 2004~\cite{Portoles2004}, Ribeiro et al. 1996~\cite{Ribeiro1996}, Schwartz et al. 1985~\cite{Schwartz1985}, Shionoiri et al. 1985~\cite{Shionoiri1985}, Shioya et al. 1992~\cite{Shioya1992}, Stage et al. 2017~\cite{Stage2017}, Sunaga et al. 1995~\cite{Sunaga1995}, Swanson et al. 1984~\cite{Swanson1984}, Tarkiainen et al. 2015~\cite{Tarkiainen2015}, Thongnopnua et al. 2005~\cite{Thongnopnua2005}, Till et al. 1984~\cite{Till1984}, Tsuruoka et al. 2007~\cite{Tsuruoka2007}, Ulm et al. 1982~\cite{Ulm1982}, VanHecken et al. 2020~\cite{VanHecken2020}, Wade et al. 1992~\cite{Wade1992}, Weisser et al. 1991~\cite{Weisser1991}, Weisser et al. 1992~\cite{Weisser1992}, and Witte et al. 1993~\cite{Witte1993}.
Tab~\ref{tab:study_table}. provides an overview of the curated studies and key information such as number of subjects, dosing protocol, route of administration, diseases, physiological impairments and genetic variants. All data have been uploaded to the pharmacokinetics database \href{https://pk-db.com}{PK-DB}~\cite{Grzegorzewski2021} and are freely accessible via the PK-DB study identifiers.
\begin{table}[h!]
\centering
\caption[Curated clinical studies]{\textbf{Overview of curated clinical studies.} po: oral, iv: intravenous}
\vspace{0.1cm}
{\fontsize{7}{9}\selectfont
    
\begin{tabular}{L{0.12\textwidth} L{0.068\textwidth} L{0.057\textwidth} L{0.035\textwidth} L{0.05\textwidth} L{0.035\textwidth} L{0.05\textwidth} L{0.035\textwidth} L{0.035\textwidth} L{0.035\textwidth} L{0.035\textwidth} L{0.035\textwidth} L{0.05\textwidth}}
\textbf{Study} & \textbf{PK-DB} & \textbf{PMID} & \textbf{route} & \textbf{substance} & \textbf{dosing} & \textbf{dose [mg]} & \textbf{healthy} & \textbf{hyper\-tension} & \textbf{renal impairment} & \textbf{hepatic impairment} & \textbf{cardiac impairment} & \textbf{CES1} \\

\hline

\rowcolor{Lightgrey}
Arafat2005 \cite{Arafat2005} & \href{https://identifiers.org/pkdb:PKDB00741}{PKDB00741} & \href{https://pubmed.ncbi.nlm.nih.gov/15985045/}{15985045} & po & enamal & single & 10 & \checkmark &  &  &  &  & \\
Baba1990 \cite{Baba1990} & \href{https://identifiers.org/pkdb:PKDB00740}{PKDB00740} & \href{https://pubmed.ncbi.nlm.nih.gov/2165799/}{2165799} & po & enamal & single & 10 & \checkmark &  &  & \checkmark &  & \\
\rowcolor{Lightgrey}
Biollaz1982 \cite{Biollaz1982} & \href{https://identifiers.org/pkdb:PKDB00768}{PKDB00768} & \href{https://pubmed.ncbi.nlm.nih.gov/6289859/}{6289859} & po & enamal & single & 10 & \checkmark &  &  &  &  & \\
Dayyih2017 \cite{Dayyih2017} & \href{https://identifiers.org/pkdb:PKDB00769}{PKDB00769} & \href{https://pubmed.ncbi.nlm.nih.gov/28894622/}{28894622} & po & enamal & single & 20 & \checkmark &  &  &  &  & \\
\rowcolor{Lightgrey}
Dickstein1987 \cite{Dickstein1987} & \href{https://identifiers.org/pkdb:PKDB00770}{PKDB00770} & \href{https://pubmed.ncbi.nlm.nih.gov/3034316/}{3034316} & po, iv & enamal, eat & single & 5, 10 & \checkmark &  &  &  & \checkmark & \\
Donnelly1990 \cite{Donnelly1990} & \href{https://identifiers.org/pkdb:PKDB00771}{PKDB00771} & \href{https://pubmed.ncbi.nlm.nih.gov/2154407/}{2154407} & po & enamal & single, multi & 20 &  & \checkmark &  &  &  & \\
\rowcolor{Lightgrey}
Fruncillo1987 \cite{Fruncillo1987} & \href{https://identifiers.org/pkdb:PKDB00773}{PKDB00773} & \href{https://pubmed.ncbi.nlm.nih.gov/3037160/}{3037160} & po & enamal & single, multi & 5 &  & \checkmark & \checkmark &  &  & \\
Ghosh2012 \cite{Ghosh2012} & \href{https://identifiers.org/pkdb:PKDB00746}{PKDB00746} & \href{https://pubmed.ncbi.nlm.nih.gov/21341376/}{21341376} & po & enamal & single & 20 & \checkmark &  &  &  &  & \\
\rowcolor{Lightgrey}
Gu2004 \cite{Gu2004} & \href{https://identifiers.org/pkdb:PKDB00747}{PKDB00747} & \href{https://pubmed.ncbi.nlm.nih.gov/15556550/}{15556550} & po & enamal & single & 10 & \checkmark &  &  &  &  & \\
He2014 \cite{He2014} & \href{https://identifiers.org/pkdb:PKDB00774}{PKDB00774} & \href{https://pubmed.ncbi.nlm.nih.gov/24721587/}{24721587} & po & enamal & single & 40 & \checkmark &  &  &  &  & \\
\rowcolor{Lightgrey}
Her2021 \cite{Her2021} & \href{https://identifiers.org/pkdb:PKDB00775}{PKDB00775} & \href{https://pubmed.ncbi.nlm.nih.gov/33963573/}{33963573} & po & enamal & multi & 2.5 & \checkmark &  &  &  &  & wt, G143E\\
Hockings1986 \cite{Hockings1986} & \href{https://identifiers.org/pkdb:PKDB00776}{PKDB00776} & \href{https://pubmed.ncbi.nlm.nih.gov/3011046/}{3011046} & po, iv & enamal, eat & single & 10 & \checkmark &  & \checkmark &  &  & \\
\rowcolor{Lightgrey}
Howes1991 \cite{Howes1991} & \href{https://identifiers.org/pkdb:PKDB00777}{PKDB00777} & \href{https://pubmed.ncbi.nlm.nih.gov/1932608/}{1932608} & po & enamal & single & 10 & \checkmark &  &  &  &  & \\
Ishizaki1988 \cite{Ishizaki1988} & \href{https://identifiers.org/pkdb:PKDB00750}{PKDB00750} & \href{https://pubmed.ncbi.nlm.nih.gov/2468049/}{2468049} & po & enamal & single & 10 & \checkmark &  &  &  &  & \\
\rowcolor{Lightgrey}
Johnston1992 \cite{Johnston1992} & \href{https://identifiers.org/pkdb:PKDB00778}{PKDB00778} & \href{https://pubmed.ncbi.nlm.nih.gov/1329474/}{1329474} & po & enamal & single, multi & 2.5 &  &  &  &  & \checkmark & \\
Kelly1986 \cite{Kelly1986} & \href{https://identifiers.org/pkdb:PKDB00779}{PKDB00779} & \href{https://pubmed.ncbi.nlm.nih.gov/3004546/}{3004546} & po & enamal & single, multi & 10 & \checkmark &  & \checkmark &  &  & \\
\rowcolor{Lightgrey}
Lee2003 \cite{Lee2003} & \href{https://identifiers.org/pkdb:PKDB00780}{PKDB00780} & \href{https://pubmed.ncbi.nlm.nih.gov/12772271/}{12772271} & po & enamal & single & 20 & \checkmark &  &  &  &  & \\
Lees1987a \cite{Lees1987a} & \href{https://identifiers.org/pkdb:PKDB00781}{PKDB00781} & \href{https://pubmed.ncbi.nlm.nih.gov/3034468/}{3034468} & po & enamal & multi & 10 & \checkmark &  & \checkmark &  &  & \\
\rowcolor{Lightgrey}
Lowenthal1985 \cite{Lowenthal1985} & \href{https://identifiers.org/pkdb:PKDB00782}{PKDB00782} & \href{https://pubmed.ncbi.nlm.nih.gov/2998676/}{2998676} & po & enamal & single & 10 & \checkmark &  & \checkmark &  &  & \\
Lu2009 \cite{Lu2009} & \href{https://identifiers.org/pkdb:PKDB00745}{PKDB00745} & \href{https://pubmed.ncbi.nlm.nih.gov/19046843/}{19046843} & po & enamal & single & 10 & \checkmark &  &  &  &  & \\
\rowcolor{Lightgrey}
MacDonald1993 \cite{MacDonald1993} & \href{https://identifiers.org/pkdb:PKDB00783}{PKDB00783} & \href{https://pubmed.ncbi.nlm.nih.gov/9114905/}{9114905} & po & enamal & single & 10 & \checkmark &  & \checkmark &  &  & \\
Marzo2002 \cite{Marzo2002} & \href{https://identifiers.org/pkdb:PKDB00784}{PKDB00784} & \href{https://pubmed.ncbi.nlm.nih.gov/12040965/}{12040965} & po & enamal & multi & 10, 20 & \checkmark &  &  &  &  & \\
\rowcolor{Lightgrey}
Matalka2002 \cite{Matalka2002} & \href{https://identifiers.org/pkdb:PKDB00785}{PKDB00785} & \href{https://pubmed.ncbi.nlm.nih.gov/12165071/}{12165071} & po & enamal & single & 10, 20 & \checkmark &  &  &  &  & \\
Moffett2014 \cite{Moffett2014} & \href{https://identifiers.org/pkdb:PKDB00786}{PKDB00786} & \href{https://pubmed.ncbi.nlm.nih.gov/27129124/}{27129124} & po & enamal & single & 10 & \checkmark &  &  &  &  & \\
\rowcolor{Lightgrey}
Mujais1992 \cite{Mujais1992} & \href{https://identifiers.org/pkdb:PKDB00787}{PKDB00787} & \href{https://pubmed.ncbi.nlm.nih.gov/1310827/}{1310827} & iv & eat & multi & 1.5 & \checkmark &  &  &  &  & \\
Najib2003a \cite{Najib2003a} & \href{https://identifiers.org/pkdb:PKDB00788}{PKDB00788} & \href{https://pubmed.ncbi.nlm.nih.gov/14520685/}{14520685} & po & enamal & single & 20 & \checkmark &  &  &  &  & \\
\rowcolor{Lightgrey}
Niopas2004 \cite{Niopas2004} & \href{https://identifiers.org/pkdb:PKDB00789}{PKDB00789} & \href{https://pubmed.ncbi.nlm.nih.gov/15112862/}{15112862} & po & enamal & single & 20 & \checkmark &  &  &  &  & \\
Oguchi1993 \cite{Oguchi1993} & \href{https://identifiers.org/pkdb:PKDB00790}{PKDB00790} & \href{https://pubmed.ncbi.nlm.nih.gov/8504625/}{8504625} & po & enamal & single & 5 & \checkmark & \checkmark & \checkmark &  &  & \\
\rowcolor{Lightgrey}
Ohnishi1989 \cite{Ohnishi1989} & \href{https://identifiers.org/pkdb:PKDB00791}{PKDB00791} & \href{https://pubmed.ncbi.nlm.nih.gov/2543535/}{2543535} & po & enamal & single & 10 & \checkmark &  &  & \checkmark &  & \\
Portoles2004 \cite{Portoles2004} & \href{https://identifiers.org/pkdb:PKDB00792}{PKDB00792} & \href{https://pubmed.ncbi.nlm.nih.gov/24936102/}{24936102} & po & enamal & single & 20 & \checkmark &  &  &  &  & \\
\rowcolor{Lightgrey}
Ribeiro1996 \cite{Ribeiro1996} & \href{https://identifiers.org/pkdb:PKDB00793}{PKDB00793} & \href{https://pubmed.ncbi.nlm.nih.gov/8839663/}{8839663} & po & enamal & single & 20 & \checkmark &  &  &  &  & \\
Schwartz1985 \cite{Schwartz1985} & \href{https://identifiers.org/pkdb:PKDB00794}{PKDB00794} & \href{https://pubmed.ncbi.nlm.nih.gov/2410720/}{2410720} & po & enamal & single & 2.5, 5, 10, 20, 40 &  & \checkmark &  &  & \checkmark & \\
\rowcolor{Lightgrey}
Shionoiri1985 \cite{Shionoiri1985} & \href{https://identifiers.org/pkdb:PKDB00749}{PKDB00749} & \href{https://pubmed.ncbi.nlm.nih.gov/2982051/}{2982051} & po & enamal & single & 10 &  & \checkmark &  &  &  & \\
Shioya1992 \cite{Shioya1992} & \href{https://identifiers.org/pkdb:PKDB00795}{PKDB00795} & \href{https://pubmed.ncbi.nlm.nih.gov/1322206/}{1322206} & po & enamal & single & 5 & \checkmark &  &  &  &  & \\
\rowcolor{Lightgrey}
Stage2017 \cite{Stage2017} & \href{https://identifiers.org/pkdb:PKDB00797}{PKDB00797} & \href{https://pubmed.ncbi.nlm.nih.gov/28639420/}{28639420} & po & enamal & single & 10 & \checkmark &  &  &  &  & wt, c4, c3, ca3, G143E, CES1A1c\\
Sunaga1995 \cite{Sunaga1995} & \href{https://identifiers.org/pkdb:PKDB00798}{PKDB00798} & \href{https://pubmed.ncbi.nlm.nih.gov/8582461/}{8582461} & po & enamal & single & 5 &  & \checkmark &  &  &  & \\
\rowcolor{Lightgrey}
Swanson1984 \cite{Swanson1984} & \href{https://identifiers.org/pkdb:PKDB00799}{PKDB00799} & \href{https://pubmed.ncbi.nlm.nih.gov/6097665/}{6097665} & po & enamal & single & 40 & \checkmark &  &  &  &  & \\
Tarkiainen2015 \cite{Tarkiainen2015} & \href{https://identifiers.org/pkdb:PKDB00800}{PKDB00800} & \href{https://pubmed.ncbi.nlm.nih.gov/25919042/}{25919042} & po & enamal & single & 10 & \checkmark &  &  &  &  & wt, G143E\\
\rowcolor{Lightgrey}
Thongnopnua2005 \cite{Thongnopnua2005} & \href{https://identifiers.org/pkdb:PKDB00801}{PKDB00801} & \href{https://pubmed.ncbi.nlm.nih.gov/15797799/}{15797799} & po & enamal & single & 5 & \checkmark &  &  &  &  & \\
Till1984 \cite{Till1984} & \href{https://identifiers.org/pkdb:PKDB00802}{PKDB00802} & \href{https://pubmed.ncbi.nlm.nih.gov/6091806/}{6091806} & po & enamal & single, multi & 10 & \checkmark &  &  &  &  & \\
\rowcolor{Lightgrey}
Ulm1982 \cite{Ulm1982} & \href{https://identifiers.org/pkdb:PKDB00804}{PKDB00804} & \href{https://pubmed.ncbi.nlm.nih.gov/6289858/}{6289858} & po & enamal & single & 9.8 & \checkmark &  &  &  &  & \\
VanHecken2020 \cite{VanHecken2020} & \href{https://identifiers.org/pkdb:PKDB00805}{PKDB00805} & \href{https://pubmed.ncbi.nlm.nih.gov/31411383/}{31411383} & po & enamal & single & 10 &  &  &  &  &  & \\
\rowcolor{Lightgrey}
Wade1992 \cite{Wade1992} & \href{https://identifiers.org/pkdb:PKDB00806}{PKDB00806} & \href{https://pubmed.ncbi.nlm.nih.gov/1312853/}{1312853} & po & enamal & single & 10 & \checkmark &  &  &  &  & \\
Weisser1991 \cite{Weisser1991} & \href{https://identifiers.org/pkdb:PKDB00807}{PKDB00807} & \href{https://pubmed.ncbi.nlm.nih.gov/1647954/}{1647954} & po & enamal & single & 10 & \checkmark &  &  &  &  & \\
\rowcolor{Lightgrey}
Weisser1992 \cite{Weisser1992} & \href{https://identifiers.org/pkdb:PKDB00808}{PKDB00808} & \href{https://pubmed.ncbi.nlm.nih.gov/1330574/}{1330574} & po & enamal & single & 10 &  & \checkmark & \checkmark &  &  & \\
Witte1993 \cite{Witte1993} & \href{https://identifiers.org/pkdb:PKDB00809}{PKDB00809} & \href{https://pubmed.ncbi.nlm.nih.gov/8394796/}{8394796} & po & enamal & single & 10 &  & \checkmark &  &  &  & \\
\end{tabular}
}
\end{table}\label{tab:study_table}
\subsection{PBPK model of enalapril}
\label{sec:results_model}
Using the curated dataset, a physiologically based pharmacokinetic (PBPK) model was developed for enalapril (Fig.~\ref{fig:enalapril_model}). The model is hierarchical, with the top layer representing the whole body, including lung, liver, kidney, intestine and finally the rest of the body. The transport of enalapril and enalaprilat is modelled via the systemic circulation. Enalapril can be given orally as enalapril maleate tablets or capsules, or intravenously as enalapril solution. Enalaprilat may be given orally or intravenously. Both compounds can be administered as a single dose or as repeated doses at regular intervals, e.g. once daily for a week. 

After oral administration, the tablet dissolves in the stomach and the drug is released into the intestine. Enalapril is absorbed in the small intestine by passive diffusion without the use of transport proteins. The absorbed drug enters the body through the portal vein. The unabsorbed fraction of enalapril in the intestine is excreted in the faeces (Fig.~\ref{fig:enalapril_model}B)~\cite{Knutter2008, Smeets2020a, Smeets2020}.

Enalapril is then taken up by the liver with the help of the transport proteins OATP1B1 and OATP1B3~\cite{Liu2006a}. The metabolism of enalapril to enalaprilat is catalysed by the enzyme carboxylesterase 1 (CES1)~\cite{Thomsen2014}. The enalaprilat formed is exported into the systemic circulation by the multi-drug resistance-associated protein 4 (MRP4)~\cite{Ferslew2014} (Fig.~\ref{fig:enalapril_model}C). 

Enalapril and enalaprilat are eliminated by glomerular filtration and tubular secretion via the kidneys (Fig.~\ref{fig:enalapril_model}D).

The model version used in this study is \href{https://zenodo.org/doi/10.5281/zenodo.11067564}{0.9.5}~\cite{enalapril-model-v0.9.5} (see Methods for licensing and repository details). The implementation of hepatic, renal and cardiac impairment and of CES1 activity is described together with their effects in Sec.~\ref{sec:results_conditions}.
% An overview of the SBML graphs of the whole body is shown in Fig.~\ref{fig:ena_sbml_body}, the intestine in Fig.~\ref{fig:ena_sbml_intestine}, the liver in Fig.~\ref{fig:ena_sbml_liver}, and the kidneys in Fig.~\ref{fig:ena_sbml_kidney}.


% \clearpage

%%% SBML graph enalapril model %%%

% \begin{figure}[!ht]
%     \centering  
%     % \includegraphics[width=1.07\textwidth]{figures/model/enalapril_body.png}
%     \caption[SBML whole-body model]{\textbf{SBML species-reaction graph of the whole-body model of enalapril.}}
%     \label{fig:ena_sbml_body}
% \end{figure}

% \begin{figure}[!ht]
%     \begin{minipage}{0.26\textwidth}
%         % \includegraphics[width=1.2\textwidth]{figures/model/enalapril_intestine.png}
%         \caption[SBML intestine model]{\textbf{SBML graph of the intestine model.}}
%         \label{fig:ena_sbml_intestine}
%     \end{minipage}
%     \hfil
%     \begin{minipage}{0.26\textwidth}
%         % \includegraphics[width=1.2\textwidth]{figures/model/enalapril_liver.png}
%         \caption[SBML liver model]{\textbf{SBML graph of the liver model.}}
%         \label{fig:ena_sbml_liver}
%     \end{minipage}
%     \hfil
%     \begin{minipage}{0.26\textwidth}
%         % \includegraphics[width=1.2\textwidth]{figures/model/enalapril_kidney.png}
%         \caption[SBML kidney model]{\textbf{SBML graph of the kidney model.}}
%         \label{fig:ena_sbml_kidney}
%     \end{minipage}
% \end{figure} 
% \clearpage


% \clearpage

\subsection{Dose dependency}
\label{sec:results_dose_dependency}
% \begin{figure}[!h]
%     \centering
%     % \includegraphics[width=0.9\linewidth]{figures//results/DoseDependencyExperiment_Fig_application_po.png}
%     \caption[Dose dependency simulation]{\textbf{Dose dependency of enalapril pharmacokinetics.} A representative simulation showing the enalapril and enalaprilat plasma concentrations and urinary and faecal amounts for varying doses of enalapril maleate namely, 2.5mg, 5mg, 10mg, 20mg, 40mg.}
%     \label{fig:DoseDependency}
% \end{figure}

To characterise the basic behaviour of the model, the dose dependency of enalapril pharmacokinetics was simulated for the most common doses of enalapril found in our literature search (Fig.~\ref{fig:DoseDependency}). These simulations use the final parameter set (Tab.~\ref{tab:Parameter_Fitting}), whose derivation is described in the next section. With increasing dose, the plasma concentrations of enalapril and enalaprilat increase, as do the amounts excreted in urine and faeces. The pharmacokinetics of enalapril is much faster than that of enalaprilat, with plasma concentrations returning to baseline after approximately 5 hours. In contrast, the elimination of enalaprilat is much slower, with larger doses being eliminated after approximately 24 hours. Enalapril and enalaprilat are excreted in the urine, with the unabsorbed fraction of enalapril being excreted in the faeces. The total amount of drug is recovered as enalapril and enalaprilat in urine and faeces.


The predicted dose dependence was compared with independent clinical data from Schwartz et al. 1985~\cite{Schwartz1985}, in which hypertensive patients (Fig.~\ref{fig:Schwartz1985_HT}) and patients with chronic heart failure (Fig.~\ref{fig:Schwartz1985_CHF}) received different doses of enalapril.


% \begin{figure}[!h]
%     \centering
%     % \includegraphics[width=0.4\linewidth]{figures//supplementary_data/Schwartz1985_Fig1_HT.png}
%     \caption[Hypertensive patients (Schwartz1985)]{\textbf{Dose dependency healthy subjects.} Timecourse data for Schwartz1985 for hypertensive patients with its respective simulations~\cite{Schwartz1985}.}
%     \label{fig:Schwartz1985_HT}
% \end{figure}

% \begin{figure}[!h]
%     \centering
%     % \includegraphics[width=0.4\linewidth]{figures//supplementary_data/Schwartz1985_Fig1_CHF.png}
%     \caption[Chronic heart failure patients (Schwartz1985)]{\textbf{Dose dependency chronic heart failure subjects.} Timecourse data for Schwartz1985 for patients having chronic heart failure with its respective simulations~\cite{Schwartz1985}.}
%     \label{fig:Schwartz1985_CHF}
% \end{figure}


\subsection{Parameter Fitting}
\label{sec:results_parameter_fitting}
{\color{red}A subset of the model parameters was estimated from plasma and urine time courses of enalapril and enalaprilat from the curated studies~\cite{Arafat2005, Ghosh2012, Ishizaki1988} (Tab.~\ref{tab:Parameter_Fitting}). Only data from healthy and hypertensive subjects after single application were used for fitting; all data after multiple applications, in disease states and with non-wild-type CES1 activity were reserved for validation. Tab.~\ref{tab:study_table} provides information on the individual study protocols, including route of administration and dose.} 
Fitting was performed in three steps to fit different subsets of parameters: i) intravenous administration of enalaprilat; ii) oral administration of enalapril; iii) metabolism of enalapril to enalaprilat.

% \FloatBarrier
\begin{table}[h!]
\centering
\small
\caption[Fitted model parameters]{\textbf{Parameters fitted in the model.}}
\label{tab:Parameter_Fitting}
\vspace{0.1cm}
\begin{tabular}{L{0.17\textwidth} L{0.350\textwidth} P{0.07\textwidth} P{0.05\textwidth} P{0.05\textwidth} P{0.105\textwidth}}
 \toprule
\textbf{Parameter} & \textbf{Description} & \textbf{Optimal value} & \textbf{Lower bound} & \textbf{Upper bound}& \textbf{Unit}\\
% \midrule
\hline
\rowcolor{Lightgrey}
\multicolumn{6}{c}{Intravenous administration of enalaprilat}\\
\hline
\texttt{KI\_\_EATEX\_k} & rate of urinary excretion of enalaprilat & 0.527922 & 0.1 & 10 & $\frac{1}{min}$ \newline \\
\texttt{ftissue\_eat} & tissue distribution of enalaprilat & 9.995439 & 0.01 & 10  & $\frac{l}{min}$ \newline \\
\texttt{K\textsubscript{p}\_eat} & partition coefficient between tissue and plasma for enalaprilat & 0.101726 & 0.01 & 1  & dimensionless \newline \\
\hline
\rowcolor{Lightgrey}
\multicolumn{6}{c}{Oral administration of enalapril}\\
\hline
\texttt{GU\_\_ENAABS\_k} & rate of absorption of enalapril in the intestine & 0.100000 & 0.1 & 100 & $\frac{1}{min}$ \\
\texttt{ftissue\_ena} & tissue distribution of enalapril & 9.999987 & 0.01 & 10 & $\frac{l}{min}$ \\
\texttt{K\textsubscript{p}\_ena} & partition coefficient between tissue and plasma for enalapril & 0.304908 & 0.1 & 10 & dimensionless \\
\texttt{LI\_\_ENAIM\_Vmax} & V\textsubscript{max} for uptake of enalapril in the liver & 0.748446 & 0.001 & 100 & $\frac{mmol\cdot l}{min}$ \\
\texttt{KI\_\_ENAEX\_k} & rate of urinary excretion of enalapril & 0.910320 & 0.001 & 10 & $\frac{1}{min}$ \\
\hline
\rowcolor{Lightgrey}
\multicolumn{6}{c}{Metabolism of enalapril to enalaprilat}\\
\hline
\texttt{LI\_\_ENA2EAT\_V\textsubscript{max}} & rate of conversion of enalapril to enalaprilat in the liver & 0.640305 & 0.001 & 100 & $\frac{mmol\cdot l}{min}$ \\
\texttt{LI\_\_EATEX\_V\textsubscript{max}} & rate of export of enalaprilat out of the liver & 0.001480 & 0.001 & 100 & $\frac{mmol\cdot l}{min}$ \\
\texttt{LI\_\_EATEX\_K\textsubscript{m}\_eat} & Michaelis constant for export of enalaprilat out of the liver & 0.430425 & 0.001 & 1 & mM \\
 \bottomrule
\end{tabular}
\end{table}
% \FloatBarrier

%EATIV
% \newpage
\subsubsection{Intravenous administration of enalaprilat}

% \FloatBarrier
% \begin{figure}[!h]
%     \centering
%     %   \includesvg[width=0.43\textwidth]{figures/fitting/datapoint_scatter_eat.svg}
%       \hfil
%     %   \includesvg[width=0.43\textwidth]{figures/fitting/residual_scatter_eat.svg}
%       \hfil
%     %   \includesvg[width=0.43\textwidth]{figures/fitting/cost_scatter_eat.svg}
%     \caption[Parameter fitting results for iv dose of enalaprilat]{\textbf{Parameter fitting results for intravenous dose of enalaprilat.} The top left panel shows the goodness of fit plot between model prediction and experimental data for the resulting optimal parameters. The top right panel shows the weighted residuals of the optimal parameter fit for the data used in the parameter fitting (table~\ref{tab:Parameter_Fitting}). The bottom panel shows the initial vs. final cost of fitting, and all three studies showed a decrease in cost for all three studies.} 
%     \label{fig:prediction_residuals_eat}
% \end{figure}
% \FloatBarrier

The first round of fitting was done for intravenous enalaprilat using only data from such subjects and only data on enalaprilat concentrations in plasma and urine. Fig.~\ref{fig:prediction_residuals_eat} shows an overview of the results of the parameter fitting consisting of a goodness-of-fit plot and residuals. Parameter fitting improved the agreement between data and model predictions, with the final model showing good agreement between experimental data and model predictions as can be seen from the cost function.

% \FloatBarrier
% \begin{figure}[!h]
%     \centering
%     %   \includesvg[width=0.9\textwidth]{figures/fitting/study fits/Dickstein1987_fm_enalaprilat_EATIV5_normal.svg}
%       \vfil
%     %   \includesvg[width=0.9\textwidth]{figures/fitting/study fits/Hockings1986_fm_EATIV10_young_enalaprilat.svg}
%       \vfil
%     %   \includesvg[width=0.9\textwidth]{figures/fitting/study fits/Mujais1992_fm_EATB_EATINF_enalaprilat.svg}
%     \caption[Fitted simulations for iv dose enalaprilat]{\textbf{Parameter fitting results for intravenous dose of enalaprilat.}} 
%     \label{fig:fitting_result_eat}
% \end{figure}
% \FloatBarrier

Fig.~\ref{fig:fitting_result_eat} shows the simulations for each of the fitted studies together with the reference data. All three studies use an intravenous dose of enalaprilat and were therefore used jointly to determine the optimal enalaprilat parameters. After fitting, the simulated time courses agree with the reference data. For Mujais et al. 1992~\cite{Mujais1992}, which also reports standard deviations, the simulations deviate slightly from individual data points but remain within the error bars.

%ENAPO
\subsubsection{Oral administration of enalapril}

% \FloatBarrier
% \begin{figure}[!h]
%     \centering
%     %   \includesvg[width=0.43\textwidth]{figures/fitting/datapoint_scatter_ena.svg}
%       \hfil
%     %   \includesvg[width=0.43\textwidth]{figures/fitting/residual_scatter_ena.svg}
%       \hfil
%     %   \includesvg[width=0.43\textwidth]{figures/fitting/cost_scatter_ena.svg}
%     \caption[Parameter fitting results for po dose of enalapril]{\textbf{Parameter fitting results for oral dose of enalapril.}} 
%     \label{fig:prediction_residuals_ena}
% \end{figure}
% \FloatBarrier

The second round of fitting was done for oral enalapril using data on enalapril concentrations in plasma and urine. Fig.~\ref{fig:prediction_residuals_ena} shows the residual plots for this fitting. Most of the reference data points are in close agreement with the predicted data. The fitting runs show a decrease in cost for most studies, with a few studies showing slightly more than optimal cost values. Fig.~\ref{fig:fitting_result_ena} shows the simulations for two representative studies with the optimal parameters. Oguchi et al. 1993~\cite{Oguchi1993} show a very good fit between the reference data and the simulated data. Even for Biollaz et al. 1982~\cite{Biollaz1982}, where the fit is not optimal, the simulated data show the same trend as the reference data and the simulation is within the standard deviation of the reference data. The parameters for the oral dose of enalapril have thus been optimised. 

% \FloatBarrier
% \begin{figure}[!h]
%     \centering
%     %   \includesvg[width=0.8\textwidth]{figures/fitting/study fits/Biollaz1982_fm_ENA_enalapril.svg}
%       \vfil
%     %   \includesvg[width=0.8\textwidth]{figures/fitting/study fits/Oguchi1993_fm_po_ENA5_1_enalapril.svg}
%       \vfil
%     %   \includesvg[width=0.75\textwidth]{figures/fitting/study fits/Oguchi1993_fm_po_ENA5_1_enalapril_urine.svg}
%     \caption[Fitted simulations for po dose enalapril]{\textbf{Parameter fitting results for oral dose of enalapril.}} 
%     \label{fig:fitting_result_ena}
% \end{figure}
% \FloatBarrier


% ENA_EAT
% \newpage
\subsubsection{Metabolism of enalapril to enalaprilat}

% \FloatBarrier
% \begin{figure}[!h]
%     \centering
%     %   \includesvg[width=0.43\textwidth]{figures/fitting/datapoint_scatter_enaeat.svg}
%       \hfil
%     %   \includesvg[width=0.43\textwidth]{figures/fitting/residual_scatter_enaeat.svg}
%       \hfil
%     %   \includesvg[width=0.43\textwidth]{figures/fitting/cost_scatter_enaeat.svg}
%     \caption[Parameter fitting results for enalapril metabolism]{\textbf{Parameter fitting results for enalapril to enalaprilat metabolism.}} 
%     \label{fig:prediction_residuals_ena_eat}
% \end{figure}
% \FloatBarrier

The final round of fitting focused on the conversion of enalapril to enalaprilat. The predicted data agree with the reference data except for a few data points (Fig.~\ref{fig:prediction_residuals_ena_eat}). This adjustment resulted in reduced costs in most cases, but also increased costs in some cases. Fig.~\ref{fig:fitting_result_enaeat} shows representative plots, Ohnishi et al. 1989~\cite{Ohnishi1989} and Kelly et al. 1986~\cite{Kelly1986}, for enalaprilat concentrations. For Ohnishi et al. 1989, the simulations are not in complete agreement with the reference data, but the general trend is followed and the simulation is within the standard deviation. The simulation of plasma and urine enalaprilat concentrations in Kelly et al. 1986 is in almost perfect agreement with the reference data, demonstrating that all rounds of fitting have contributed to a significant improvement in the model. 

% \FloatBarrier
% \begin{figure}[!h]
%     \centering
%     %   \includesvg[width=0.75\textwidth]{figures/fitting/study fits/Ohnishi1989_fm_ENA10_enalaprilat_healthy.svg}
%       \vfil
%     %   \includesvg[width=0.75\textwidth]{figures/fitting/study fits/Ohnishi1989_fm_enalaprilat_cumulative amount_healthy.svg}
%       \vfil
%     %   \includesvg[width=0.75\textwidth]{figures/fitting/study fits/Kelly1986_fm_po_S1_enalaprilat.svg}
%       \vfil
%     %   \includesvg[width=0.75\textwidth]{figures/fitting/study fits/Kelly1986_fm_po_S1_enalaprilat_urine.svg}
%     \caption[Fitted simulations for enalapril metabolism]{\textbf{Parameter fitting results for enalapril to enalaprilat metabolism.}} 
%     \label{fig:fitting_result_enaeat}
% \end{figure}
% \FloatBarrier

In summary, the parameters have been optimised, resulting in much better agreement between model predictions and clinical data. The optimal parameters (Tab.~\ref{tab:Parameter_Fitting}) have been incorporated into the model version \href{https://zenodo.org/doi/10.5281/zenodo.11067564}{0.9.5}~\cite{enalapril-model-v0.9.5}.


\subsection{Simulation of curated clinical studies}
\label{sec:results_simulations}
With the fitted parameters, simulations were carried out for all datasets of the curated studies and compared with the reference data reported in each study.

\subsubsection{Pharmacokinetics}
\label{sec:results_pk_simulations}

The respective graphs are reported in the Supplementary Data as 
Arafat2005 Fig.~\ref{fig:Supp_Arafat2005} \cite{Arafat2005},
Baba1990 Fig.~\ref{fig:Supp_Baba1990} \cite{Baba1990},
Biollaz1982 Fig.~\ref{fig:Supp_Biollaz1982} \cite{Biollaz1982},
Dayyih2017 Fig.~\ref{fig:Supp_Dayyih2017} \cite{Dayyih2017},
Dickstain1987 Fig.~\ref{fig:Supp_Dickstein1987}, \ref{fig:Supp_Dickstein1987_CHF}, \ref{fig:Supp_Dickstein1987_IV}, \ref{fig:Supp_Dickstein1987_IV_CHF}, \ref{fig:Supp_Dickstein1987_IV_EAT}, \ref{fig:Supp_Dickstein1987_IV_EAT_CHF} \cite{Dickstein1987},
Donelly1990 Fig.~\ref{fig:Supp_Donnelly1990} \cite{Donnelly1990},
Fruncillo1987 Fig.~\ref{fig:Supp_Fruncillo1987} \cite{Fruncillo1987}, 
Ghosh2012 Fig.~\ref{fig:Supp_Ghosh2012} \cite{Ghosh2012},
Gu2004 Fig.~\ref{fig:Supp_Gu2004} \cite{Gu2004},
He2014 Fig.~\ref{fig:Supp_He2024} \cite{He2014},
Her2021 Fig.~\ref{fig:Supp_Her2021} \cite{Her2021}, 
Hockings1986 Fig.~\ref{fig:Supp_Hockings1986} \cite{Hockings1986},
Howes1981 Fig.~\ref{fig:Supp_Howes1991} \cite{Howes1991},
Ishizaki1988 Fig.~\ref{fig:Supp_Ishizaki1988} \cite{Ishizaki1988},
Johnston1992 Fig.~\ref{fig:Supp_Johnston1992} \cite{Johnston1992},
Kelly1986 Fig.~\ref{fig:Supp_Kelly1986} \cite{Kelly1986},
Lee2023 Fig.~\ref{fig:Supp_Lee2023} \cite{Lee2003},
Lees\-1987a Fig.~\ref{fig:Supp_Lees1987a} \cite{Lees1987a},
Lowenthal1985 Fig.~\ref{fig:Supp_Lowenthal1985} \cite{Lowenthal1985},
Lu2009 Fig.~\ref{fig:Supp_Lu2009} \cite{Lu2009},
MacDonald1993 Fig.~\ref{fig:Supp_MacDonald1993} \cite{MacDonald1993},
Marzo2002 Fig.~\ref{fig:Supp_Marzo2002} \cite{Marzo2002},
Matalka2002 Fig.~\ref{fig:Supp_Matalka2002} \cite{Matalka2002},
Moffett2014 Fig.~\ref{fig:Supp_Moffett2014} \cite{Moffett2014},
Mujais1992 Fig.~\ref{fig:Supp_Mujais1992} \cite{Mujais1992},
Najib2003a Fig.~\ref{fig:Supp_Najib2003a} \cite{Najib2003a},
Niopas2004 Fig.~\ref{fig:Supp_Niopas2011} \cite{Niopas2004},
Oguchi1993 Fig.~\ref{fig:Supp_Oguchi1993} \cite{Oguchi1993},
Ohnishi1989 Fig.~\ref{fig:Supp_Ohnishi1989} \cite{Ohnishi1989},
Portales2004 Fig.~\ref{fig:Supp_Portoles2004} \cite{Portoles2004},
Ribeiro1996 Fig.~\ref{fig:Supp_Ribeiro1996} \cite{Ribeiro1996},
Schwartz1985 Fig.~\ref{fig:Supp_Schwartz1985}, \ref{fig:Supp_Schwartz1985_CHF} \cite{Schwartz1985},
Shionoiri1985 Fig.~\ref{fig:Supp_Shionoiri1985} \cite{Shionoiri1985}, 
Shioya1992 Fig.~\ref{fig:Supp_Shioya1992} \cite{Shioya1992},
Stage\-2017 Fig.~\ref{fig:Supp_Stage2017_GAIN}, \ref{fig:Supp_Stage2017_LOSS} \cite{Stage2017},
Sunaga1995 Fig.~\ref{fig:Supp_Sunaga1995} \cite{Sunaga1995}, 
Swanson1984 Fig.~\ref{fig:Supp_Swanson1984} \cite{Swanson1984}, 
Tarkiainen2015 Fig.~\ref{fig:Supp_Tarkiainen2015} \cite{Tarkiainen2015},
Thonnopnua2005 Fig.~\ref{fig:Supp_Thongnopnua2005}~\cite{Thongnopnua2005},
Till1984 Fig.~\ref{fig:Supp_Till1984} \cite{Till1984},
Tsuruoaka2007 Fig.~\ref{fig:Supp_Tsuruoka2007} \cite{Tsuruoka2007},
Ulm1982 Fig.~\ref{fig:Supp_Ulm1982} \cite{Ulm1982},
VanHecken2020 Fig.~\ref{fig:Supp_VanHecken2020} \cite{VanHecken2020},
Wade1992 Fig.~\ref{fig:Supp_Wade1992} \cite{Wade1992},
Weisser1991 Fig.~\ref{fig:Supp_Weisser1991} \cite{Weisser1991},
Weisser1992 Fig.~\ref{fig:Supp_Weisser1992} \cite{Weisser1992},
Witte1993 Fig.~\ref{fig:Supp_Witte1993} \cite{Witte1993}.

Two example simulations are shown below Fig.~\ref{fig:Biollaz1982} and Fig.~\ref{fig:Dayyih2017}.

% \begin{figure}[!h] 
%     \centering
%     % \includegraphics[width=0.47\textwidth]{figures/supplementary_data/Biollaz1982_Fig1.png}
%     \caption[Biollaz1982 healthy 10mg enalapril]{\textbf{Example simulation Biollaz1982.} Biollaz1982 studied the pharmacokinetics of enalapril/-at in healthy, normotensive subjects with 10mg administration of enalapril maleate~\cite{Biollaz1982}.}
%     \label{fig:Biollaz1982}
% \end{figure}

% \begin{figure}[!h] 
%     \centering 
%     % \includegraphics[width=0.47\textwidth]{figures/supplementary_data/Dayyih2017_Fig5&6.png}
%     \caption[Dayyih2017 healthy 20mg enalapril]{\textbf{Example simulation Dayyih2017.} Dayyih2017 compared the pharmacokinetics of enalapril/-at in healthy subjects with 20mg administration of two different formulations of enalapril maleate~\cite{Dayyih2017}.}
%     \label{fig:Dayyih2017}
% \end{figure}

Biollaz et al. 1982~\cite{Biollaz1982} (Fig.~\ref{fig:Biollaz1982}) studied the pharmacokinetics of enalapril in 12 normotensive subjects aged between 22 and 33 years. They analysed the plasma concentrations of enalapril and enalaprilat after a single dose of 10 mg enalapril maleate. Our model was able to simulate the time courses of enalapril and enalaprilat accurately. Dayyih et al.2017~\cite{Dayyih2017} (Fig.~\ref{fig:Dayyih2017}) took a slightly different approach and attempted to compare the pharmacokinetics between two different formulations of enalapril maleate. Their results showed that there was no significant difference in the pharmacokinetics of the two formulations. Our model was also able to simulate a similar time course that is in good agreement with the data for both formulations.

\subsubsection{Pharmacodynamics}
\label{sec:results_pd_simulations}


%%% Specific conditions %%%
\subsection{Effect of specific conditions on pharmacokinetics}
\label{sec:results_conditions}
The fitted model was used to investigate how hepatic impairment, renal impairment and CES1 activity affect the pharmacokinetics of enalapril and enalaprilat. This section first describes how these factors are implemented in the model and then, for each factor in turn, presents the corresponding results. For each factor, parameter scans of a single oral dose of 20~mg enalapril (see Methods) were combined with simulations of the curated studies in the corresponding patient populations. None of these disease or genotype data were used for parameter fitting, so the comparisons with the reference data are independent model validations.

\subsubsection{Implementation of organ impairment and CES1 activity}
% \begin{figure}[!h]
%     \centering
%     % \includegraphics[width=1.07\linewidth]{figures/model/hepatic_renal_impairment.png}
%     \caption[Hepatic and renal impairment]{\textbf{Overview of hepatic and renal impairment.} \textbf{A)} Liver impairment was modelled as a gradual increase in cirrhosis by scaling liver function with the parameter \texttt{f\_cirrhosis} from 0.0 (no cirrhosis) to 0.95 (critical cirrhosis). Normal liver function: black (1.0), mild cirrhosis: light blue (0.40), moderate cirrhosis: blue (0.7) and severe cirrhosis: dark blue (0.81), corresponding to Child Pugh Turette classes CPT A, CPT B and CPT C respectively. Cirrhosis is modelled as a combination of reduction in functional liver volume and shunting of blood around the liver, both of which result in reduced liver function.
%     \textbf{B)} Renal impairment has been modelled as a progressive decline in renal function by scaling all renal processes with the factor \texttt{f\_renal\_function}, where 1.0 represents normal renal function and 0.0 represents no renal function. Normal renal function: black (1.0), mild renal impairment: light green (0.69), moderate renal impairment: green (0.32) and severe renal impairment: dark green (0.19).
%     }
%     \label{fig:renal_hepatic}
% \end{figure}

Renal and hepatic impairment have far-reaching effects on the pharmacokinetics of drugs. The model was therefore adapted to simulate hepatic, renal and combined hepato-renal impairment, as well as cardiac impairment and CES1 activity, each implemented as a scaling factor on the affected model processes (Fig.~\ref{fig:renal_hepatic}). The figure provides an overview of the different degrees of hepatic and renal impairment considered in our model. 

Accounting for hepatic impairment in a pharmacokinetic model requires quantification of the degree of hepatic impairment. Indocyanine green is a compound that is routinely used in clinical settings to test liver function in subjects. Köller et al 2021 established a pharmacokinetic model of indocyanine green in patients with hepatic impairment in two studies~\cite{Koller2021, Koller2021a}. Thus, the ICG levels reported in the enalapril studies~\cite{Baba1989, Ohnishi1989} helped us to determine the respective degrees of hepatic impairment.

Similarly, in renal impairment, creatinine clearance or glomerular filtration rate are reliable parameters to assess renal function and hence the degree of renal impairment~\cite{Hockings1986}. The reported values of these renal markers were used to determine the degree of renal impairment.

Cardiac impairment was represented by a 30\% reduction of cardiac output (\texttt{f\_cardiac\_output} = 0.70). CES1 activity was modelled by scaling the CES1-catalysed conversion of enalapril to enalaprilat (\texttt{f\_ces1}): 1.0 for wild type, 0.1 for the loss-of-function variant G143E, 1.2 for the active allele and 1.0 for CES1A1c, with genotypes carrying two different alleles represented by the mean of both (e.g., 0.55 for WT/G143E). The resulting PBPK model allows the pharmacokinetics of enalapril and enalaprilat to be simulated under various degrees of hepatic, renal and cardiac impairment and for different CES1 genotypes.

%%% Hepatic Impairment %%%
% \clearpage
% \newpage
\subsubsection{Hepatic impairment}
\label{sec:results_hepatic_scans}

We used the model to investigate the effect of hepatic impairment on enalapril pharmacokinetics. In addition, simulations were performed for the two curated studies investigating enalapril pharmacokinetics in cirrhotic patients. The simulated data agreed with the reference data for the general trend of rise and fall of enalapril and enalaprilat concentrations.

% \begin{figure}[!h]
%     \centering
%     % \includegraphics[width=1.0\linewidth]{figures/scan_data/EnalaprilParameterScan_fig_ena_po__renal__hepatic_scan__timecourse.png}
%     \caption[Timecourse for hepatic function scan]{\textbf{Timecourse for hepatic function scan.} Time courses for different degrees of cirrhosis. Cirrhosis in red, normal function in black. The scan was repeated for different degrees of renal function (control: black, mild renal impairment: light green, moderate renal impairment: green, and severe renal impairment: dark green.}
%     \label{fig:timecourse_renal_hepatic}
% \end{figure}

Time courses were simulated for enalapril pharmacokinetics as shown in Fig.~\ref{fig:timecourse_renal_hepatic}. Plasma enalapril concentrations increased with increasing hepatic and renal impairment. Decreased efficiency of the liver to convert enalapril to enalaprilat and decreased ability of the kidney to excrete enalapril from the blood lead to increased concentrations. The opposite effect is observed with enalaprilat, where dramatic increases in plasma concentrations are observed. With enalapril, urinary concentrations increase with progressive stages of cirrhosis, but decrease with increasing renal impairment. With enalaprilat, urine concentrations decrease with increasing cirrhosis.    

% \begin{figure}[!h]
%     \centering
%     % \includegraphics[width=1.0\linewidth]{figures/scan_data/EnalaprilParameterScan_fig_ena_po__enalapril__renal__hepatic_scan__parameters.png}
%     \caption[Pharmacokinetics of enalapril for hepatic function scan]{\textbf{Pharmacokinetic parameters of enalapril for renal function scan.} Pharmacokinetic parameters of enalapril for varying degrees of hepatic function. The scan was repeated for different degrees of renal function (control: black, mild renal impairment: light green, moderate renal impairment: green, and severe renal impairment: dark green.}
%     \label{fig:parameter_ena_renal_hepatic}
% \end{figure}
% \begin{figure}[!h]
%     \centering
%     % \includegraphics[width=1.0\linewidth]{figures/scan_data/EnalaprilParameterScan_fig_ena_po__enalaprilat__renal__hepatic_scan__parameters.png}
%     \caption[Pharmacokinetics of enalaprilat for hepatic function scan]{\textbf{Pharmacokinetic parameters of enalaprilat for renal function scan.}
%     Pharmacokinetic parameters of enalaprilat for varying degrees of hepatic function. The scan was repeated for different degrees of renal function (control: black, mild renal impairment: light green, moderate renal impairment: green, and severe renal impairment: dark green.}
%     \label{fig:parameter_eat_renal_hepatic}
% \end{figure}

For the pharmacokinetic parameters of enalapril, AUC and t\textsubscript{1/2} increased with increasing degree of cirrhosis, whereas hepatic clearance decreased. No effect of degree of cirrhosis on renal clearance was observed. The AUC for enalaprilat decreased slightly with increasing cirrhosis, but showed an increase with increasing renal impairment. The t\textsubscript{1/2} hardly changed in these scenarios. The renal clearance of enalaprilat decreased significantly with increasing cirrhosis and renal impairment.

The model was used to simulate the pharmacokinetics of enalapril and enalaprilat in patients with varying degrees of hepatic impairment~\cite{Baba1990, Ohnishi1989}. The figures below are the graphs for the simulations of the two studies that focused on patients with hepatic impairment or patients with different degrees of cirrhosis.

% \begin{figure}[!h]
%     \centering
%     % \includegraphics[width=0.7\linewidth]{figures//supplementary_data/Baba1990_Fig1.png}
%     \caption[Hepatic impairment Baba1990]{\textbf{Simulation of hepatic impairment for Baba1990.} Normal hepatic function: black, cirrhosis: blue.~\cite{Baba1990}}
%     \label{fig:Baba1990}
% \end{figure}

% \begin{figure}[!h]
%     \centering
%     % \includegraphics[width=1.08\linewidth]{figures//supplementary_data/Ohnishi1989_Fig1.png}
%     \caption[Hepatic impairment Ohnishi1989]{\textbf{Simulation of hepatic impairment for Ohnishi1989.} Normal hepatic function: black, cirrhosis: blue.~\cite{Ohnishi1989}}
%     \label{fig:Ohnishi1989}
% \end{figure}


%%% Renal Impairment %%%
% \clearpage
% \newpage
\subsubsection{Renal impairment}
\label{sec:results_renal_scans}

We used the model to investigate the effect of renal impairment on the pharmacokinetics of enalapril. As shown in Fig.~\ref{fig:timecourse_hepatic_renal}, the model reproduced the plasma time courses of enalapril and enalaprilat across the modeled conditions of hepatic impairment. Plasma concentrations of enalapril and enalaprilat decreased more rapidly with increasing renal activity, i.e., less renal impairment, whereas urinary recovery of both substances increased with higher renal activity. In the absence of cirrhosis, enalapril levels fell more quickly because its metabolism to enalaprilat proceeded efficiently. In severe cirrhosis, by contrast, enalapril persisted longer in the blood and reached a higher peak concentration than in the control group, while enalaprilat concentrations failed to reach the maxima observed in the control group. Under cirrhotic conditions, urinary enalapril was substantially higher and urinary enalaprilat substantially lower than in the control group.

% \begin{figure}[!h]
%     \centering
%     % \includegraphics[width=1.0\linewidth]{figures//scan_data/EnalaprilParameterScan_fig_ena_po__hepatic__renal_scan__timecourse.png}
%     \caption[Timecourse for renal function scan]{\textbf{Timecourse for renal function scan.} Time courses for different degrees of renal function. Decreased renal function in red, increased renal function in blue, normal function in black. The scan was repeated for different degrees of cirrhosis (control: black, mild cirrhosis: light blue (CPT A), moderate cirrhosis: blue (CPT B) and severe cirrhosis: dark blue (CPT C).}
%     \label{fig:timecourse_hepatic_renal}
% \end{figure}

For the pharmacokinetic parameters (Fig.~\ref{fig:parameter_ena_hepatic_renal}, Fig.~\ref{fig:parameter_eat_hepatic_renal}), cirrhotic conditions also increased the AUC and decreased the elimination rate (t\textsubscript{1/2}) of enalapril, consistent with enalapril remaining longer in plasma when hepatic conversion to enalaprilat is impaired. The AUC and t\textsubscript{1/2} differed between cirrhotic groups under reduced renal activity, but not at baseline or higher renal activity. Renal and hepatic clearance of enalapril increased with increasing renal function, while hepatic clearance decreased with progressive stages of cirrhosis, consistent with reduced hepatic metabolic capacity. For enalaprilat, the AUC and t\textsubscript{1/2} decreased steadily with increasing renal function and with increasing degree of cirrhosis.

% \begin{figure}[!h]
%     \centering
%     % \includegraphics[width=1.0\linewidth]{figures//scan_data/EnalaprilParameterScan_fig_ena_po__enalapril__hepatic__renal_scan__parameters.png}
%     \caption[Pharmacokinetics of enalapril for renal function scan]{\textbf{Pharmacokinetic parameters of enalapril for renal function scan.} Pharmacokinetic parameters of enalapril for varying degrees of renal function. The scan was repeated for different degrees of cirrhosis (control: black, mild cirrhosis: light blue (CPT A), moderate cirrhosis: blue (CPT B) and severe cirrhosis: dark blue (CPT C).}
%     \label{fig:parameter_ena_hepatic_renal}
% \end{figure}

% \begin{figure}[!h]
%     \centering
%     % \includegraphics[width=0.8\linewidth]{figures//scan_data/EnalaprilParameterScan_fig_ena_po__enalaprilat__hepatic__renal_scan__parameters.png}
%    \caption[Pharmacokinetics of enalaprilat for renal function scan]{\textbf{Pharmacokinetic parameters of enalaprilat for renal function scan.} Pharmacokinetic parameters of enalaprilat for varying degrees of renal function. The scan was repeated for different degrees of cirrhosis (control: black, mild cirrhosis: light blue (CPT A), moderate cirrhosis: blue (CPT B) and severe cirrhosis: dark blue (CPT C).}
%     \label{fig:parameter_eat_hepatic_renal}
% \end{figure}

The model was used to simulate the pharmacokinetics of enalapril and enalaprilat in patients with varying degrees of renal impairment~\cite{Fruncillo1987, Hockings1986, Kelly1986, Lees1987a, Lowenthal1985, Weisser1992}. The simulated data showed varying degrees of goodness of fit to the reference data and the fit improved with more iterations of the parameterisation. The figures below show plots of the simulated and reference data for the renal impairment studies.

% \begin{figure}[!h]
%     \centering
%     % \includegraphics[width=0.9\linewidth]{figures//supplementary_data/Fruncillo1987_Fig7.png}
%     \caption[Renal impairment Fruncillo1987]{\textbf{Simulation of renal impairment for Fruncillo1987.} Normal renal function: black, renal impairment: green.~\cite{Fruncillo1987}}
%     \label{fig:Fruncillo1987}
% \end{figure}
% \begin{figure}[!h]
%     \centering
%     % \includegraphics[width=0.7\linewidth]{figures//supplementary_data/Hockings1986_Fig_EATIV10.png}
%     \caption[Renal impairment Hockings1986 IV]{\textbf{Simulation of renal impairment for Hockings1986.} Normal renal function: black, renal impairment: green.~\cite{Hockings1986}}
%     \label{fig:Hockings1986_IV}
% \end{figure}
% \begin{figure}[!h]
%     \centering
%     % \includegraphics[width=0.8\linewidth]{figures//supplementary_data/Hockings1986_Fig_ENAPO10.png}
%     \caption[Renal impairment Hockings1986 PO]{\textbf{Simulation of renal impairment for Hockings1986.} Normal renal function: black, renal impairment: green.~\cite{Hockings1986}}
%     \label{fig:Hockings1986_PO}
% \end{figure}
% \begin{figure}[!h]
%     \centering
%     % \includegraphics[width=0.8\linewidth]{figures/supplementary_data/Kelly1986_Fig_multi.png}
%     \caption[Renal impairment Kelly1986 single]{\textbf{Simulation of renal impairment for Kelly1986.} Normal renal function: black, renal impairment: green.~\cite{Kelly1986}}
%     \label{fig:Kelly1986_single}
% \end{figure}
% \begin{figure}[!h]
%     \centering
%     % \includegraphics[width=0.8\linewidth]{figures/supplementary_data/Kelly1986_Fig_single.png}
%     \caption[Renal impairment Kelly1986 multi]{\textbf{Simulation of renal impairment for Kelly1986.} Normal renal function: black, renal impairment: green.~\cite{Kelly1986}}
%     \label{fig:Kelly1986_multi}
% \end{figure}
% \begin{figure}[!h]
%     \centering
%     % \includegraphics[width=0.6\linewidth]{figures/supplementary_data/Lees1987a_Fig2.png}
%     \caption[Renal impairment Lees1987a]{\textbf{Simulation of renal impairment for Lees1987a.} Normal renal function: black, renal impairment: green.~\cite{Lees1987a}}
%     \label{fig:Lees1987a}
% \end{figure}
% \begin{figure}[!h]
%     \centering
%     % \includegraphics[width=0.7\linewidth]{figures/supplementary_data/Lowenthal1985_Fig1.png}
%     \caption[Renal impairment Lowenthal1985]{\textbf{Simulation of renal impairment for Lowenthal1985.} Normal renal function: black, renal impairment: green.~\cite{Lowenthal1985}}
%     \label{fig:Lowenthal1985}
% \end{figure}
% \begin{figure}[!h]
%     \centering
%     % \includegraphics[width=1.0\linewidth]{figures/supplementary_data/Weisser1992_Fig1.png}
%     \caption[Renal impairment Weisser1992]{\textbf{Simulation of renal impairment for Weisser1992.} Normal renal function: black, renal impairment: green.~\cite{Weisser1992}}
%     \label{fig:Weisser1992}
% \end{figure}

%%% CES1 Genotype %%%
% \clearpage
% \newpage
\subsubsection{CES1 genotype}
\label{sec:results_ces1_scans}

The model was also used to investigate the effect of CES1 genotype on enalapril pharmacokinetics. Simulations were performed for a range of CES1 activity against 4 degrees of cirrhosis. \ref{fig:timecourse_hepatic_CES1} It was observed that CES1 and hepatic impairment have similar and synergistic effects on the pharmacokinetics of enalapril. In both cases, the conversion of enalapril to enalaprilat is inhibited. Therefore, under non-mutated conditions, enalapril concentrations tend to be lower than in the case of reduced CES1 activity and higher than in the case of increased CES1 activity. On the other hand, enalaprilat concentrations tend to be higher under conditions of higher CES1 activity and they tend to be lower under conditions of lower CES1 activity. Urinary enalapril concentrations tend to increase when CES1 activity decreases and vice versa. The opposite trend is observed for urinary enalaprilat concentrations. The degree of cirrhosis does not seem to have any effect other than to increase the total enalapril concentration and decrease that of enalaprilat in all the above cases.

% \begin{figure}[!h]
%     \centering
%     % \includegraphics[width=1.0\linewidth]{figures/scan_data/EnalaprilParameterScan_fig_ena_po__hepatic__CES1__timecourse.png}
%     \caption[Timecourse for CES1 function scan (hepatic impairment)]{\textbf{Timecourse for CES1 function scan.} Time courses for different degrees of CES1 function. Decreased CES1 function in red, increased CES1 function in blue, normal function in black. The scan was repeated for different degrees of cirrhosis (control: black, mild cirrhosis: light blue (CPT A), moderate cirrhosis: blue (CPT B) and severe cirrhosis: dark blue (CPT C).}
%     \label{fig:timecourse_hepatic_CES1}
% \end{figure}

Pharmacokinetic parameters were also estimated with the simulations using our model. The AUC and total clearance of enalapril were unaffected between the 4 degrees of cirrhosis at subnormal CES1 activity, but at higher CES1 activity the AUC was increased and total clearance decreased compared to normal conditions.

% \begin{figure}[!h]
%     \centering
%     % \includegraphics[width=1.0\linewidth]{figures/scan_data/EnalaprilParameterScan_fig_ena_po__enalapril__hepatic__CES1__parameters.png}
%     \caption[Pharmacokinetics of enalapril for CES1 function scan (hepatic impairment)]{\textbf{Pharmacokinetic parameters of enalapril for CES1 function scan.} Pharmacokinetic parameters of enalapril for varying degrees of CES1 function. The scan was repeated for different degrees of cirrhosis (control: black, mild cirrhosis: light blue (CPT A), moderate cirrhosis: blue (CPT B) and severe cirrhosis: dark blue (CPT C).}
%     \label{fig:parameter_ena_hepatic_CES1}
% \end{figure}

The pharmacokinetic parameters for enalapril and enalaprilat were estimated as shown in Fig.~\ref{fig:parameter_ena_hepatic_CES1} and Fig.~~\ref{fig:parameter_eat_hepatic_CES1}. The AUC and t\textsubscript{1/2} for enalapril are seen to decrease with higher levels of CES1 activity. Hepatic clearance also decreases, but no effect on renal clearance is observed. For enalaprilat, the AUC increases with increasing CES1 activity, but the t\textsubscript{1/2} remains the same; renal clearance also increases.

% \begin{figure}[!h]
%     \centering
%     % \includegraphics[width=0.8\linewidth]{figures/scan_data/EnalaprilParameterScan_fig_ena_po__enalaprilat__hepatic__CES1__parameters.png}
%     \caption[Pharmacokinetics of enalaprilat for CES1 function scan (hepatic impairment)]{\textbf{Pharmacokinetic parameters of enalaprilat for CES1 function scan.} Pharmacokinetic parameters of enalaprilat for varying degrees of CES1 function. The scan was repeated for different degrees of cirrhosis (control: black, mild cirrhosis: light blue (CPT A), moderate cirrhosis: blue (CPT B) and severe cirrhosis: dark blue (CPT C).}
%     \label{fig:parameter_eat_hepatic_CES1}
% \end{figure}

CES1 activity was also investigated in relation to renal impairment. The plasma concentrations of enalapril and enalaprilat were dramatically higher in patients with renal impairment than in healthy patients. This is because renal impairment impairs renal filtration and reduces the excretion of both substances, resulting in a higher presence of these substances in the blood/plasma. In terms of CES1 activity, enalapril levels decreased as CES1 activity increased. Conversely, enalaprilat levels increased as CES1 activity increased.

Renal impairment only increased the concentration of enalapril and decreased that of enalaprilat. The urinary concentrations of ena\-lapril and enalaprilat showed opposite trends. Thus, a combined effect of higher CES1 activity and severe renal impairment was that the urinary concentration of enalapril decreased dramatically, whereas the urinary concentration of enalaprilat did not decrease, but was certainly excreted more rapidly. On the other hand, with reduced CES1 activity and no renal impairment, the urinary concentrations of enalapril and enalaprilat increased and decreased, respectively. Only a change in CES1 activity was associated with a change in enalapril-at concentration and the trend was the same, i.e. as above, enalapril concentration increased with higher CES1 activity and enalaprilat concentration increased with higher CES1 activity.

% \begin{figure}[!h]
%     \centering
%     % \includegraphics[width=1.0\linewidth]{figures/scan_data/EnalaprilParameterScan_fig_ena_po__renal__CES1__timecourse.png}
%     \caption[Timecourse renal with CES1]{Timecourses for enalapril with varying degrees of renal impairment against CES1 activity}
%     \caption[Timecourse for CES1 function scan (renal impairment)]{\textbf{Timecourse for CES1 function scan.} Time courses for different degrees of CES1 function. Decreased CES1 function in blue, increased CES1 function in red, normal function in black. The scan was repeated for different degrees of renal function (control: black, mild renal impairment: light green, moderate renal impairment: green, and severe renal impairment: dark green.}
%     \label{fig:timecourse_renal_CES1}
% \end{figure}

% \begin{figure}[!h]
%     \centering
%     % \includegraphics[width=1.0\linewidth]{figures/scan_data/EnalaprilParameterScan_fig_ena_po__enalapril__renal__CES1__parameters.png}
%     \caption[Pharmacokinetics of enalapril for CES1 function scan (renal impairment)]{\textbf{Pharmacokinetic parameters of enalapril for CES1 function scan.} Pharmacokinetic parameters of enalapril for varying degrees of CES1 function. The scan was repeated for different degrees of renal function (control: black, mild renal impairment: light green, moderate renal impairment: green, and severe renal impairment: dark green.}
%     \label{fig:parameter_ena_renal_CES1}
% \end{figure}

Pharmacokinetic parameters were also estimated for CES1 activity as shown in Fig.~\ref{fig:parameter_ena_renal_CES1} and Fig.~\ref{fig:parameter_eat_renal_CES1}. For enalapril, AUC and t\textsubscript{1/2} both increase with decreasing CES1 activity, but remain the same at higher than normal CES1 activity. In addition, the greater the degree of renal impairment, the higher the AUC at lower CES1 activity; the same is true for t\textsubscript{1/2}. Hepatic clearance of enalapril has been shown to increase with increasing CES1 activity, while renal clearance remains the same. Only renal impairment has been shown to affect the renal clearance of enalapril.

For enalaprilat, the AUC increases with CES1 activity, but this change is more apparent with higher degrees of renal impairment. The t\textsubscript{1/2} remains the same for different CES activities and only increases with higher degrees of renal impairment. The renal clearance of enalaprilat is interesting because at lower levels of CES1 activity it is lower with less renal impairment; however, beyond normal CES1 activity, the renal clearance of enalaprilat decreases with higher degrees of renal impairment.

% \begin{figure}[!h]
%     \centering
%     % \includegraphics[width=0.8\linewidth]{figures/scan_data/EnalaprilParameterScan_fig_ena_po__enalaprilat__renal__CES1__parameters.png}
%     \caption[Pharmacokinetics of enalaprilat for CES1 function scan (renal impairment)]{\textbf{Pharmacokinetic parameters of enalaprilat for CES1 function scan.} Pharmacokinetic parameters of enalaprilat for varying degrees of CES1 function. The scan was repeated for different degrees of renal function (control: black, mild renal impairment: light green, moderate renal impairment: green, and severe renal impairment: dark green.}
%     \label{fig:parameter_eat_renal_CES1}
% \end{figure}

The model was used to simulate the pharmacokinetics of enalapril and enalaprilat in patients with different degrees of CES1 function~\cite{Her2021, Stage2017, Tarkiainen2015}. Simulations were also performed with a special focus on studies comparing the pharmacokinetics between variants of the CES1 gene, which gives rise to the CES1 enzyme responsible for the conversion of enalapril to enalaprilat (see Fig.~\ref{fig:Her2021_ces1}, Fig.~\ref{fig:Stage2017_ces1_GAIN}, Fig.~\ref{fig:Stage2017_ces1_LOSS} and Fig.~\ref{fig:Tarkiainen2015_ces1}).
% \begin{figure}[!h]
%     \centering
%     % \includegraphics[width=0.5\linewidth]{figures/supplementary_data/Her2021_Fig3.png}
%     \caption[CES1 Her2021]{\textbf{Simulation for CES1 activity Her2021.} Normal CES1 function: black, increased CES1 function: blue, decreased CES1 function: red.~\cite{Her2021}}
%     \label{fig:Her2021_ces1}
% \end{figure}
% \begin{figure}[!h]
%     \centering
%     % \includegraphics[width=0.4\linewidth]{figures/supplementary_data/Stage2017_Fig_GAIN.png}
%     \caption[CES1 Stage2017]{\textbf{Simulation for CES1 activity Stage2017.} Normal CES1 function: black, increased CES1 function: blue, decreased CES1 function: red.~\cite{Stage2017}}
%     % \includegraphics[width=0.4\linewidth]{figures/supplementary_data/Stage2017_Fig_LOSS.png}
%     \caption[CES1 Stage2017]{\textbf{Simulation for CES1 activity Stage2017.} Normal CES1 function: black, increased CES1 function: blue, decreased CES1 function: red.~\cite{Stage2017}}
%     \label{fig:Stage2017_ces1}
% \end{figure}
% \clearpage
% \begin{figure}[!h]
%     \centering
%     % \includegraphics[width=0.7\linewidth]{figures/supplementary_data/Tarkiainen2015_Fig1.png}
%     \caption[CES1 Tarkiainen2015]{\textbf{Simulation for CES1 activity Tarkiainen2015.} Normal CES1 function: black, increased CES1 function: blue, decreased CES1 function: red.~\cite{Tarkiainen2015}}
%     \label{fig:Tarkiainen2015_ces1}
% \end{figure}

% \clearpage
% \newpage
\subsection{Summary}
Together, these results show that the model reproduces the dose-dependent pharmacokinetics of enalapril and enalaprilat in plasma and urine and, after stepwise parameter optimisation on single-dose data from healthy and hypertensive subjects, describes a broad set of independent clinical studies. Hepatic impairment shifted the balance from enalaprilat formation towards enalapril persistence, increasing enalapril AUC and t\textsubscript{1/2} while lowering urinary enalaprilat. Renal impairment raised plasma concentrations of both compounds and reduced their urinary recovery. Reduced CES1 activity affected the conversion step in the same direction as hepatic impairment, whereas increased CES1 activity had the opposite effect.
